\section{Kết luận và định hướng tương lai}
\label{sec:conclusion}

Nghiên cứu này đã thiết kế, cài đặt từ đầu bằng NumPy và đánh giá thực nghiệm có kiểm soát một quy trình hai giai đoạn kết hợp giữa lọc đốm không gian, phân đoạn ngưỡng toàn cục và phân loại học máy theo vùng nhằm phát hiện vết dầu loang trên ảnh SAR đơn kênh. Bằng việc tính toán mô-men cục bộ của bộ lọc Lee MMSE trong thời gian $\mathcal{O}(HW)$ qua bảng tích phân và huấn luyện bộ phân loại LinearSVM tối ưu bằng thuật toán Nesterov Accelerated Gradient (NAG) trên không gian 8 đặc trưng vùng (cùng hệ số chặn bias tạo thành vector tham số 9 chiều), công trình đã cung cấp một cơ sở thực nghiệm minh bạch về cả hiệu quả giảm báo nhầm lẫn các giới hạn cấu trúc của phương pháp cổ điển trên bộ dữ liệu SOS.

Kết quả kiểm thử độc lập trên 839 ảnh Sentinel-1A (chứa $3\,182$ đối tượng dầu tham chiếu) và 776 ảnh ngoài cảm biến ALOS PALSAR ghi nhận các chỉ số định lượng chính xác như sau:
\begin{itemize}
    \item \textbf{Triệt tiêu cảnh báo sai ở giai đoạn phân loại vùng}: Khi so sánh có kiểm soát ở cùng cấu hình tạo ứng viên (Lee $7 \times 7, L=4$ + Otsu + $A_{\min} = 64$), bộ phân loại LinearSVM đã loại bỏ $76.6\%$ số vùng báo nhầm (cắt giảm $9\,572$ vùng tối giả, đưa tỷ lệ cảnh báo sai từ $14.89$ xuống $3.48$ FP/ảnh), đặc biệt giảm $99.2\%$ báo nhầm trên các ảnh biển sạch $R_0$ (từ $24.32$ xuống $0.19$ FP/ảnh, khoảng tin cậy Bootstrap 95\% $[0.06, 0.35]$, với $25/31$ ảnh đạt $0$ FP).
    \item \textbf{Cải thiện độ chính xác cấp đối tượng và cấp điểm ảnh}: LinearSVM nâng Precision từ $0.0753$ lên $0.2541$ (tăng $3.37$ lần) và nâng Object F1 từ $0.1220 \ [0.1124, 0.1328]$ lên $0.2805 \ [0.2652, 0.2980]$ ($\Delta\text{F1} = +0.1586 \ [+0.1474, +0.1706]$). Bộ phân loại giữ lại $97.8\%$ số đối tượng đúng do Giai đoạn 1 đề xuất ($996/1\,018$ TP), đạt độ thu hồi tổng thể $\text{Recall} = 0.3130$ trên toàn bộ $3\,182$ đối tượng Ground Truth. Về mặt diện tích điểm ảnh, hệ thống đạt Global Pixel IoU $= 0.6614$ trên toàn tập và Macro Pixel IoU trung bình từng ảnh đạt $0.577 \pm 0.261$ ($n = 814$).
    \item \textbf{Đánh giá ngoài cảm biến trên ALOS PALSAR}: Khi áp dụng trực tiếp mô hình đã khóa sang 776 ảnh ALOS PALSAR băng L mà không huấn luyện lại, quy trình đạt Object F1 $= 0.3724$, Global Pixel IoU $= 0.6423$ và tỷ lệ báo nhầm $2.97$ FP/ảnh ($0.07$ FP/ảnh trên 45 ảnh biển sạch $R_0$), phản ánh khả năng tương thích tốt trên các cảnh sự cố tràn dầu có độ tương phản cao ở chế độ diện tích trung bình $R_2$.
    \item \textbf{Hiệu năng tính toán}: Quy trình thực thi hoàn toàn trên CPU mất trung bình $50.2 \pm 14.7\text{ ms/ảnh}$ kích thước $256 \times 256$ (thông lượng đo trên 4 tiến trình song song), nhanh hơn phương pháp phân ngưỡng trực tiếp trên ảnh thô nhờ giảm hơn $72\%$ số thành phần liên thông cần xử lý.
\end{itemize}

Quan trọng hơn, phân tích đối chứng và giải phẫu lỗi xuyên suốt bài báo đã làm sáng tỏ bốn phát hiện khoa học cốt lõi về mặt phương pháp luận:
\begin{enumerate}
    \item \textit{Hiện tượng suy biến của trọng số Lee MMSE trên ảnh SAR nén 8-bit}: Trên các bộ dữ liệu đã qua tiền xử lý đa nhìn và lượng tử hóa 8-bit như SOS, hệ số biến thiên cục bộ $C_I^2 = v/m^2$ tại phần lớn các vùng đồng nhất đều nhỏ hơn nhiều so với ngưỡng phương sai nhiễu lý thuyết $1/L$ khi chọn $L \in \{1, 4\} \ll \text{ENL}_{\text{eff}}$. Điều này khiến trọng số $W(x, y) = \max(0, 1 - \frac{1/L}{C_I^2})$ bằng $0$ trên diện rộng và làm bộ lọc Lee $7 \times 7$ hoạt động tương đương thống kê với bộ lọc trung bình hộp Mean $7 \times 7$ ($\Delta\text{F1} = +0.0014$, khoảng tin cậy 95\% $[-0.0003, +0.0033]$ chứa 0).
    \item \textit{Sự lệch pha cấu trúc giữa nhãn huấn luyện Purity và thước đo Hungarian IoU}: Bộ phân loại vùng đạt F1 tới $80.74\%$ khi đánh giá theo nhãn tinh khiết ($\mathrm{Purity} \ge 0.50$), nhưng hơn một nửa số dự đoán dương tính bị tính là False Positive ở cấp đối tượng ($>1\,250 / 2\,215$ FP trên tập Validation) thực chất là các mảnh dầu thật bị phân ngưỡng Otsu cắt vụn khiến từng mảnh riêng lẻ không đạt ngưỡng ghép cặp 1--1 $\IoU \ge 0.50$ với mặt nạ GT nguyên khối.
    \item \textit{Nút cổ chai độ thu hồi và suy biến ước lượng nền ở chế độ ngập dầu}: Trần Recall của toàn hệ thống bị khóa bởi bước phân ngưỡng Otsu toàn cục ($\le 0.3199$, trong đó các vết dầu nhỏ nội vùng chỉ đạt \texttt{Interior Recall} $= 10.48\%$). Đồng thời, trên các patch ngập dầu ($R_3, R_4$), việc chỉ lấy các điểm ảnh ngoài ứng viên làm nền khiến ước lượng $\mu_{\text{bg}}$ bị thiên lệch về phía sáng và làm suy giảm hiệu lực của đặc trưng tương phản nền.
    \item \textit{Đa cộng tuyến đặc trưng và hạn chế cấp độ patch}: Đặc trưng \texttt{boundary\_gradient} ($f_6$) có tương quan đơn biến dương với nhãn dầu ($r = +0.3632$) nhưng nhận trọng số âm ($w_6 = -0.639$) trong mô hình LinearSVM do đóng vai trò biến ức chế (suppressor variable) trước sự cộng tuyến mạnh với độ lệch chuẩn nội vùng ($r = +0.8504$) và tương phản nền ($r = +0.6716$). Ngoài ra, việc dữ liệu gốc thiếu siêu dữ liệu cảnh chụp (scene ID) và độ phân giải vật lý khiến đánh giá hiện tại giới hạn ở cấp độ patch $256 \times 256$.
\end{enumerate}

Từ những kết luận thực nghiệm trên, các hướng phát triển tiếp theo được định hình rõ ràng:
\begin{enumerate}
    \item \textbf{Hợp nhất vùng hậu phân loại và đồng bộ hóa hàm mục tiêu}: Bổ sung bước gom cụm không gian (spatial region merging / morphological bridging) đối với các ứng viên lân cận cùng được dự đoán là dầu nhằm khắc phục hiện tượng phân mảnh đối tượng và thu hẹp khoảng cách giữa nhãn huấn luyện Purity với thước đo đánh giá Hungarian IoU.
    \item \textbf{Ước lượng tự động $\text{ENL}_{\text{eff}}$ và cổng kiểm định đa đỉnh (Bimodality Gate)}: Tự động ước lượng số lượt nhìn hiệu dụng từ vùng biển đồng nhất để hiệu chỉnh tham số $L$ cho bộ lọc Lee trên ảnh 8-bit, đồng thời tích hợp kiểm định thống kê đa đỉnh trước bước phân ngưỡng nhằm ngăn chặn việc ép cắt Otsu trên các khung ảnh biển sạch ($R_0$) hoặc ngập trọn dầu ($R_4$).
    \item \textbf{Mở rộng quy mô cửa sổ lọc và chuẩn hóa theo siêu dữ liệu cảnh gốc}: Khảo sát mở rộng kích thước cửa sổ lọc ($9 \times 9, 11 \times 11$) kết hợp với các mô hình đề xuất ứng viên đa tỷ lệ (như CFAR thích nghi hoặc mạng phân đoạn nhẹ) trên dữ liệu Sentinel-1 Level-1 GRD gốc có đầy đủ tọa độ địa lý để đánh giá độc lập theo từng cảnh chụp (scene-level split).
\end{enumerate}
